\section*{第 \ref{ch:g6:fractions} 章 --- 分数初步}

\begin{solution}{exo:g6:fractions:1}
蛋糕：吃掉 $\frac38$。\quad 巧克力：$\frac{7}{12}$。\quad
派：$\frac66 = 1$（整个派）。
\end{solution}

\begin{solution}{exo:g6:fractions:2}
五份中涂了色的有四份。$\frac45 < 1$（五份中的四份
少于整条）；$\frac65 > 1$（六分之五是一整条
再多出五分之一）。
\end{solution}

\begin{solution}{exo:g6:fractions:3}
在按三等分标注的直线上：$\frac13$ 在 $0$ 之后一格；
$\frac53$ 在 $2$ 之前一格（因为 $2 = \frac63$）；$\frac73$ 在
$2$ 之后一格。所以 $\frac53$ 与 $\frac73$ 都
相距 $\frac13$，分居 $2$ 的两侧。
\end{solution}

\begin{solution}{exo:g6:fractions:4}
$\frac94 = 2 + \frac14$；\quad
$\frac{13}{5} = 2 + \frac35$；\quad
$\frac{12}{6} = 2$（恰好）；\quad
$\frac{25}{8} = 3 + \frac18$。
\end{solution}

\begin{solution}{exo:g6:fractions:5}
$\frac23 = \frac{8}{12}$（乘 $4$）；\quad
$\frac{15}{20} = \frac34$（除以 $5$）；\quad
$\frac{7}{10} = \frac{21}{30}$（乘 $3$）；\quad
$\frac56 = \frac{10}{12}$（除以 $2$）。
\end{solution}

\begin{solution}{exo:g6:fractions:6}
$\frac68 = \frac34$；\quad
$\frac{15}{25} = \frac35$；\quad
$\frac{18}{24} = \frac{9}{12} = \frac34$；\quad
$\frac{35}{7} = 5$（整数）。
\end{solution}

\begin{solution}{exo:g6:fractions:7}
$\frac12$ 乘 $86$：$86 \div 2 = 43$。

$\frac34$ 乘 $60$：$60 \div 4 = 15$，再 $15 \times 3 = 45$。

$\frac25$ 乘 $35$：$35 \div 5 = 7$，再 $7 \times 2 = 14$。

$\frac{7}{10}$ 乘 $250$：$250 \div 10 = 25$，再 $25 \times 7 = 175$。
\end{solution}

\begin{solution}{exo:g6:fractions:8}
$\frac12 > \frac13$：二分之一的每一份比三分之一的每一份大。
$\frac35 > \frac25$：份的大小相同，份数更多。
$\frac34 = \frac68 < \frac78$：八分之七比八分之六大。
\end{solution}

\begin{solution}{exo:g6:fractions:9}
$\frac12 = 0.5$；\quad $\frac34 = 0.75$；\quad
$\frac{7}{10} = 0.7$；\quad $\frac94 = 2.25$。与众不同的是
$\frac13 = 0.333\dots$，它的小数写法永不结束，因此它不是
小数。
\end{solution}

\begin{solution}{exo:g6:fractions:10}
坐公交车的人：$\frac34$ 乘 $28$：$28 \div 4 = 7$，再
$7 \times 3 = 21$ 名学生。步行的人：$28 - 21 = 7$ 名学生（即
剩下的四分之一：$28 \div 4 = 7$）。检验：$21 + 7 = 28$。
\end{solution}

\begin{solution}{exo:g6:fractions:11}
积蓄的三分之二是 $18$ 欧元，因此\emph{一}份，即三分之一，是
$18 \div 2 = 9$ 欧元，而全部积蓄（三份）是
$9 \times 3 = 27$ 欧元。检验：$\frac23$ 乘 $27$ 是 $18$。
\end{solution}

\begin{solution}{exo:g6:fractions:12}
用掉的水占水箱的比例从 $\frac58$ 变为 $\frac48$（二分之一即
八分之四）：$10$ 升对应容积的
$\frac58 - \frac48 = \frac18$。因此满箱容量为
$10 \times 8 = 80$ 升。检验：$\frac58$ 乘 $80$ 是 $50$，减去
$10$ 剩 $40$，恰好是 $80$ 的一半。
\end{solution}

\begin{solution}{pb:g6:fractions:1}
\textbf{1.} $16$ 个小方格的一半是 $8$ 个小方格。写成分数：
长条的 $\frac12$；又因为每一半是 $8$ 个小方格，
（共 $16$ 个），也可写成 $\frac{8}{16}$ --- 同一个分数的两种写法
（见\cref{prop:g6:fractions:equal}）。

\textbf{2.} 剩下的 $8$ 个小方格的一半是 $4$ 个小方格，
即原长条的 $\frac{4}{16}$。约分 $4$：
$\frac{4}{16} = \frac14$。在图上：把剩下的
一半再对半切开，得到原长条的两个四分之一 --- 一半的
一半是四分之一。

\textbf{3.} 第三口：$4$ 个小方格的一半是 $2$ 个小方格，即
长条的 $\frac{2}{16} = \frac18$。第四口：$1$ 个小方格，
即 $\frac{1}{16}$。从一口到下一口，分母
\emph{翻倍}：$2$、$4$、$8$、$16$。

\textbf{4.} 吃掉的：$8 + 4 + 2 + 1 = 15$ 个小方格（共 $16$ 个），
即长条的 $\frac{15}{16}$。剩下的：$1$ 个小方格，$\frac{1}{16}$。

\textbf{5.} $\frac12 = \frac{8}{16}$（分子分母同乘
$8$），$\frac14 = \frac{4}{16}$（乘 $4$），$\frac18 = \frac{2}{16}$
（乘 $2$），而 $\frac{1}{16}$ 保持不变。数一数十六分之几：
$8 + 4 + 2 + 1 = 15$ 份，所以总数是 $\frac{15}{16}$。

\textbf{6.} 咬下的份量：$32$、$16$、$8$、$4$、$2$、$1$ 个小方格。吃掉的
总数：$32 + 16 + 8 + 4 + 2 + 1 = 63$ 个小方格，所以剩下 $1$ 个小方格，
即长条的 $\frac{1}{64}$。

\textbf{7.} 每咬一口后，剩下的部分依次为
\[
\frac12,\quad \frac14,\quad \frac18,\quad \frac{1}{16},\quad
\frac{1}{32},\quad \frac{1}{64} :
\]
剩余部分每次减半，分母翻倍。假想的
第七口咬下后，将剩下 $\frac{1}{128}$。

\textbf{8.} 每一口只吃掉剩余部分的\emph{一半}，因此
剩余的另一半始终还在：无论咬了多少口，
总有一些剩下。因此被吃掉的总量
始终小于 $1$ --- 正如
\cref{pb:g6:decimals:1} 中由二十个 9 组成的小数，越来越接近 $3$，
却永远达不到它。

\textbf{9.} 两个分数的分子都是 $1$，而切成
$128$ 份所得的每份比切成 $100$ 份时更小，于是
$\frac{1}{128} < \frac{1}{100}$。第七口咬下之后，
剩下的部分为 $\frac{1}{128}$，因此被吃掉的部分超过
$1 - \frac{1}{100} = \frac{99}{100}$：挑战在
第七口时即告成功。

\textbf{10.} 以十六分之一为单位，各次总量分别为 $\frac{8}{16}$、
$\frac{12}{16}$、$\frac{14}{16}$、$\frac{15}{16}$。每一次新的总量
恰好落在上一次总量与 $1$ 的\emph{正中间}：
这一阶梯不断把它距离 $1$ 的剩余长度减半，
却永远不会落在这个数上。

\textbf{11.} 正确之处在于：在芝诺描述的每个阶段，
确实仍有一段距离剩下（第 8 题） --- 阶段的列举
永远不会结束。陷阱在于：各阶段在距离
\emph{和}跑步时间上都极短；跑者在各阶段上所花的
时间并不相等。把这段跑步描述为无穷多个
不断缩小的片段，并不会使跑步本身变得无穷无尽：跑者
终将到达墙边，而第 9 题表明，这一描述
会经过墙前的任意一点。

\textbf{12.} 将 $8$ 根棒各切成两半，得到 $16$ 个半块。将 $4$ 根
棒各切成四等份，得到 $16$ 个四分之一块。将 $2$ 根棒各切成
八等份，得到 $16$ 个八分之一块。将最后一根棒切成十六等份：
共 $16$ 个小块。这样一共用掉 $8 + 4 + 2 + 1 = 15$ 根棒，每个
孩子都得到每种大小的一块。

\textbf{13.} 每个孩子所得的份额为
\[
\frac12 + \frac14 + \frac18 + \frac{1}{16} = \frac{15}{16}
\]
一根棒（第 5 题）。十六份这样的份额合计
$16 \times \frac{15}{16} = 15$ 根棒
（见\cref{thm:g6:fractions:quotient}）：恰好是切开的全部，
没有剩余 --- 仅靠减半，就把 $15$ 根棒公平地分给 $16$
个孩子。

\textbf{14.} 邻居们每人各得 $\frac78$ 根棒。
分母化为 $16$：$\frac78 = \frac{14}{16}$，而我们的孩子
每人得到 $\frac{15}{16}$。因此，这 $16$ 个孩子中的每一个都多得：
$\frac{15}{16} > \frac{14}{16}$，多出十六分之一根棒。

\textbf{15.} 在正方形图中，每一口阴影部分占去
尚未涂色角落的一半，而这一未涂色的角落不断缩小，
每咬一口后，它都变为原来的一半，且
永远不会消失。下面两个事实是确定的：（i）经过任意有限多口之后，
总会剩下一部分，被吃掉的总量
始终严格小于 $1$；（ii）剩下的那部分会变得
比人们所指定的任何分数都小（例如 $\frac{1}{100}$、
$\frac{1}{1000}$、\dots）。只要咬得足够久，给出
“无穷和等于 $1$”这句话的确切含义，
这正是高中分册中\emph{极限}的任务。
\end{solution}
